\chapter{Conclusion}
\label{ch:conclusion}

I set out to find how much of the tactical picture that professional clubs buy can be recovered from one broadcast camera, and to show it in a form a coach would use. Tactical Eye does this for short clips and for full halves. A user uploads a video, follows or cancels the analysis, and then sees every player labelled by position, the referees marked, paths, team shapes and heatmaps painted on the grass in sync with the video, next to a tactical board, the formations of both teams and an AI-written report. On 20 hand-annotated frames it found 98.3\,\% of the people with no false detections and placed them on the pitch with a median error of 0.59\,m, and 119 automated tests run on every change.

The biggest surprise was how little the 30-second test clips told me. On them the pipeline looked finished. The first full half broke that impression within an hour: coaches on the touchline became players, spare balls became the ball, the referee joined a team, the tracker handed out 1,423 identities, and the movement paths jumped across the screen. None of these was a bug in the usual sense; the design had simply never met a stoppage or a 47-minute camera pan. I had also expected the neural networks to be the hard part, but they were the most reliable component. The difficult problems were in the geometry and bookkeeping around them: a camera model that wobbles, a decoder that cannot seek, a queue that hands out a finished job again after a laptop sleeps.

If I started again, I would test on a full half in the first month, and I would write the API and system tests together with the first endpoints instead of months later. I would also design the results format for long videos from the start, so that a full match could be loaded in parts.

Most of the work ran on a laptop with 8\,GB of RAM and no usable GPU, which shaped many decisions: the queue, the cancel button, the rented GPUs (under USD 2 in total for all experiments) and the protection against analysing a job twice. With footage of the whole pitch, player identities and event detection, Tactical Eye could come close to the tools professional clubs use; even now it gives an amateur coach a view of their team's shape that until recently required a stadium tracking system.
